% ============================================================
% The Collision Fluctuation Sum
% ============================================================

\documentclass[12pt]{amsart}

\usepackage{amsmath,amssymb,amsthm}
\usepackage{booktabs}
\usepackage{microtype}
\usepackage{url}

\newtheorem{theorem}{Theorem}
\newtheorem{lemma}[theorem]{Lemma}
\newtheorem{proposition}[theorem]{Proposition}
\newtheorem{corollary}[theorem]{Corollary}
\newtheorem{conjecture}[theorem]{Conjecture}
\theoremstyle{definition}
\newtheorem{definition}[theorem]{Definition}
\theoremstyle{remark}
\newtheorem{remark}[theorem]{Remark}
\newtheorem{observation}[theorem]{Observation}

\title{The Collision Fluctuation Sum}

\author{Alexander S.\ Petty}

\email{alexander.petty@gmail.com}
\date{October 2024}
\subjclass[2020]{11A63, 11N05, 11M41}

\begin{document}
\begin{abstract}
For a prime $p > b + 1$ not dividing the base~$b$, the
\emph{collision fluctuation} at lag~$\ell$ is
$\Delta_p(\ell) = C(b^{\ell} \bmod p) - \overline{C}_p$,
the deviation of the collision count from its mean over
constructive multipliers. The partial sum
\[
\Phi_x(s, \ell) = \sum_{\substack{b+1 < p \le x \\ p \nmid b}}
\frac{\Delta_p(\ell)}{p^s}
\]
converges absolutely for $\operatorname{Re}(s) > 2$.
Computation to $10^7$ suggests Mertens-rate divergence at
$s = 1$: $\Phi_x(1, \ell) \sim -\mu_{\ell} \cdot \log\log x$,
where $\mu_{\ell}$ depends on the lag and base
($\mu_1 \approx 0.68$ in base~$10$).

The divergence is controlled by the decomposition
$\Delta_p = S - QR/(p{-}1{-}b)$, where $S = C - Q$ is an
integer and $QR/(p{-}1{-}b)$ is a deterministic
correction from the spectral class $R = (p{-}1) \bmod b$.
The bilateral symmetry of the repetend forces $C(g)$ to be
even for all $g$ and constrains $S$ to have the same parity
as $Q$.
\end{abstract}
\maketitle

% ============================================================
\section{Introduction}

This paper studies the fluctuation of the collision count across primes.

The preceding paper~\cite{paper12} decomposed the fractional
field Dirichlet series into a mean part
(controlled by the mean formula from the gate width
paper~\cite{paper11}) and a fluctuation part. The mean part
converges for $\operatorname{Re}(s) > 2$. The present paper
studies the fluctuation part.

The classical prime number theorem is equivalent to the
statement that $\sum \mu(n)/n$ converges, where $\mu$ is the
M\"obius function~\cite{davenport,hardy}. This convergence
encodes the cancellation of multiplicative structure across the
integers at the edge of the critical strip.

Computation to $10^7$ suggests that the collision fluctuation
sum diverges at $s = 1$, growing like $\log\log x$ in a
Mertens-type growth law~\cite{mertens}. The divergence arises
from a persistent negative bias in the collision deviation: the
specific multiplier $b^{\ell}$ is systematically below the mean
collision count. A later paper investigates whether a centered
version of the fluctuation sum converges at $s = 1$.

% ============================================================
\section{The Fluctuation Sum}

\begin{definition}
For a prime $p > b + 1$ not dividing $b$, the \emph{collision
fluctuation} at lag~$\ell$ is
\[
\Delta_p(\ell) = C(b^{\ell} \bmod p) - \overline{C}_p,
\]
where $\overline{C}_p$ is the mean collision count over
constructive nonidentity multipliers (those $g \ne 1$ with
$C(g) > 0$). The \emph{fluctuation partial sum} is
\[
\Phi_x(s, \ell) = \sum_{\substack{b+1 < p \le x \\ p \nmid b}}
\frac{\Delta_p(\ell)}{p^s}.
\]
The digit-partitioning prime $p = b + 1$ (when prime) is
excluded because $\overline{C}_p$ is undefined there
($p - 1 - b = 0$).
\end{definition}

\begin{proposition}[Analytic convergence]\label{prop:convergence}
The series $\Phi(s, \ell) = \lim_{x \to \infty} \Phi_x(s, \ell)$
converges absolutely for $\operatorname{Re}(s) > 2$.
\end{proposition}

\begin{proof}
The fluctuation satisfies $|\Delta_p(\ell)| \le p - 1$
trivially. Thus $|\Delta_p(\ell)|/p^s = O(p^{1 - s})$.
The sum $\sum_p p^{1 - s}$ converges for
$\operatorname{Re}(s) > 2$.
\end{proof}

\begin{remark}
The variance scaling~\cite{paper12} gives
$\operatorname{Var}(\Delta) = O(Q)$ across multipliers
at fixed~$p$, suggesting typical $|\Delta_p| = O(\sqrt{Q})
= O(\sqrt{p/b})$. If this bound held pointwise for the
specific multiplier $b^{\ell}$, convergence would extend to
$\operatorname{Re}(s) > 3/2$. However, the variance scaling
is a statement about the distribution across multipliers,
not a pointwise bound for a fixed lag. For a specific
lag~$\ell$, the magnitude $|\Delta_p(\ell)|$ could be as
large as~$Q$.
\end{remark}

% ============================================================
\section{The Exact Decomposition}

\begin{theorem}[Collision decomposition]\label{thm:decomp}
The collision deviation satisfies
\[
\Delta_p = S - \frac{QR}{p - 1 - b},
\]
where $S = C(b^{\ell} \bmod p) - Q$ is an integer,
$Q = \lfloor (p{-}1)/b \rfloor$, and $R = (p{-}1) \bmod b$.
\end{theorem}

\begin{proof}
The mean collision count is
$\overline{C}_p = Q(b(Q{-}1) + 2R) / (b(Q{-}1) + R)$.
Since $b(Q{-}1) + R = p - 1 - b$:
$\overline{C}_p = Q + QR/(p{-}1{-}b)$.
Therefore $\Delta_p = C - \overline{C}_p
= (C - Q) - QR/(p{-}1{-}b) = S - QR/(p{-}1{-}b)$.
\end{proof}

\begin{remark}
Computation in base~$10$ at lag~$1$ shows that $S$ takes
values in $\{-9, -7, -4, -3, -2, -1, 0, 1, 2, 3, 6, 8\}$
for primes up to $10^7$, with $S \in \{-1, 0\}$ for
approximately $60\%$ of primes. Whether $S$ is uniformly
bounded as $p \to \infty$ is not established; the trivial
bound is $|S| \le Q$.
\end{remark}

\begin{theorem}[Bilateral parity]\label{thm:parity}
For any prime $p > b$ and any $g \in \mathbb{F}_p^{\times}$,
the collision count $C(g)$ is even.
\end{theorem}

\begin{proof}
If $r$ contributes to $C(g)$ (i.e., $\delta(r) = \delta(gr)$),
then $p - r$ also contributes: $\delta(p{-}r) = b{-}1{-}\delta(r)$
and $\delta(g(p{-}r)) = \delta(-gr) = b{-}1{-}\delta(gr)$,
so $\delta(p{-}r) = \delta(g(p{-}r))$. Since $r \ne p - r$
for $r \ne p/2$ (impossible for prime $p > 2$), the
contributions pair.
\end{proof}

\begin{corollary}
$S = C - Q$ has the same parity as $Q$.
\end{corollary}

% ============================================================
\section{The Mertens Growth Law}

Computation to $10^7$ (664,574 primes) indicates that
the partial sum $\Phi_x(1, 1)$ diverges at $s = 1$, growing
like $\log\log x$:

\begin{table}[h]
\centering
\begin{tabular}{rrrr}
\toprule
primes & $p_{\max}$ & $\Phi_x(1,1)$ & $\Phi_x/\log\log p$ \\
\midrule
$1{,}000$ & $7{,}951$ & $-1.37$ & $-0.623$ \\
$10{,}000$ & $104{,}779$ & $-1.60$ & $-0.652$ \\
$100{,}000$ & $1{,}299{,}811$ & $-1.77$ & $-0.671$ \\
$664{,}574$ & $9{,}999{,}991$ & $-1.90$ & $-0.682$ \\
\bottomrule
\end{tabular}
\caption{$\Phi_x(1,1)/\log\log p$ drifts steadily toward a
constant $\mu_1 \approx 0.68$.}
\end{table}

At smaller scales (up to 8,000 primes), the growth is invisible
because $\log\log x$ grows extremely slowly:
$\log\log(10^4) \approx 2.25$ versus
$\log\log(10^7) \approx 2.78$.

\begin{conjecture}[Mertens growth]\label{conj:mertens}
The partial sum satisfies
\[
\Phi_x(1, \ell) = -\mu_{\ell} \cdot \log\log x + O(1)
\]
for a constant $\mu_{\ell} > 0$ depending on the lag and base.
\end{conjecture}

\begin{observation}\label{obs:transition}
The transition from convergence to divergence is visible
in the partial sums:

\begin{table}[h]
\centering
\begin{tabular}{rrrrrr}
\toprule
primes & $s = 1.5$ & $s = 1.2$ & $s = 1.0$ & $s = 0.9$ &
$s = 0.5$ \\
\midrule
$100$ & $-0.17$ & $-0.47$ & $-1.07$ & $-1.42$ & $-6.97$ \\
$500$ & $-0.18$ & $-0.55$ & $-1.28$ & $-2.02$ & $-17.0$ \\
$1000$ & $-0.18$ & $-0.56$ & $-1.37$ & $-2.17$ & $-21.4$ \\
$2000$ & $-0.18$ & $-0.57$ & $-1.45$ & $-2.30$ & $-26.0$ \\
\bottomrule
\end{tabular}
\caption{$\Phi_x(s, 1)$ for various $s$. For $s > 1$, partial
sums appear to stabilize. At $s = 1$, logarithmic growth.
For $s < 1$, power-law growth.}
\end{table}
\end{observation}

The divergence arises because $\Delta_p$ has a persistent
negative mean: the mean of $\Delta_p$ over primes
$p \le x$ is approximately $-0.91$ in base~$10$ at $\ell = 1$,
with $70\%$ of primes having $\Delta_p < 0$, $20\%$ having
$\Delta_p > 0$, and $10\%$ having $\Delta_p = 0$. The
bilateral symmetry constrains $C(g)$ to be even but does not
force $\Delta_p$ to have zero mean.

% ============================================================
\section{The Per-Class Structure}

The collision deviation $\Delta_p$ decomposes by spectral
class $a = p \bmod b$.

\begin{observation}\label{obs:class-means}
For each coprime class~$a$, the mean of $\Delta_p$ over
primes $p \equiv a \pmod{b}$ with $p \le 10^7$ stabilizes
to a class-specific value.

\begin{table}[h]
\centering
\begin{tabular}{rrrr}
\toprule
$p \bmod 10$ & $R$ & mean $S$ & mean $\overline{\Delta}$ \\
\midrule
$1$ & $0$ & $-1.70$ & $-1.70$ \\
$3$ & $2$ & $-0.91$ & $-1.11$ \\
$7$ & $6$ & $-0.08$ & $-0.68$ \\
$9$ & $8$ & $+0.68$ & $-0.12$ \\
\bottomrule
\end{tabular}
\caption{Per-class means for all primes up to $10^7$
in base~$10$ at $\ell = 1$. All classes have negative mean
$\overline{\Delta}$.}
\label{tab:class-means}
\end{table}
\end{observation}

The per-class collision bias for the spectral class $R = 6$
(digit~$7$, the negative pole) is
\[
\overline{\Delta}_{R=6} \approx -\frac{2}{3}.
\]
The value $2/3$ is the alignment limit of the prime~$3$ from
the first paper~\cite{paper1}. The same constant appears as
the collision bias of the negative polarity class.

The per-class mean of $S$ in the complementary class $R = 8$
(digit~$9$, the neutral pole) is $\overline{S}_{R=8} \approx
+2/3$. The bilateral symmetry exchanges these roles: the
positive $S$ in class~$9$ is absorbed by the spectral class
correction $R/b = 0.8$, leaving a small negative
$\overline{\Delta}$.

% ============================================================
\section{The Mertens Constant}

The following heuristic relates the overall Mertens constant
to the per-class biases.

By Mertens' theorem~\cite{mertens} for arithmetic
progressions~\cite{davenport}, each coprime class contributes
equally to the prime sum at leading order. The Mertens constant
is therefore approximately
\[
\mu_{\ell} \approx \frac{1}{\varphi(b)}
\sum_{\substack{a \bmod b \\ \gcd(a,b)=1}}
|\overline{\Delta}_a|,
\]
where $\varphi$ is Euler's totient function.

In base~$10$ at $\ell = 1$: the average of the four class
biases ($1.70$, $1.11$, $0.68$, $0.12$) is $0.90$. The
observed ratio $\Phi_x(1,1)/\log\log x \approx 0.68$ is
smaller; the discrepancy suggests that the absolute-value
average overstates the effective constant, since the per-class
contributions partially cancel in the signed sum. A precise
derivation of $\mu_{\ell}$ from the class means remains open.

The constant $2/3$ from~\cite{paper1} appears not as the
overall Mertens constant but as the per-class collision bias
for the negative polarity class ($R = 6$, digit~$7$). The
Mertens constant is a weighted average in which $2/3$ is one
of four terms.

% ============================================================
\section{Structure of the Cancellation}

The partial cancellation in $\Phi_x(1, \ell)$ arises from the
sawtooth structure of the collision count established in the
preceding paper~\cite{paper12}.

At each prime $p$, the collision count decomposes as
$C(b^{\ell} \bmod p) = Q + S$, where $S$ is a signed sum of
$Q$ terms, each $\pm 1$. The sign of each term is determined by
a comparison of two sawtooth sequences modulo~$p$: one
controlled by the collision parameter $c = b(1 - b^{\ell})^{-1}
\bmod p$, the other by the base~$b$.

As $p$ ranges over the primes, the collision parameter $c$
traces a path through the residues modulo~$p$. The
equidistribution of this path, controlled by the M\"obius
deranging family, determines how many of the $Q$ signed terms
are positive versus negative. The persistent negative bias means
that $c$ systematically occupies positions with more negative
than positive sawtooth terms.

% ============================================================
\section{Remarks}

\subsection*{The boundary at $s = 1$}

The Riemann Hypothesis is equivalent to $\sum \mu(n)/n^s$
converging for all $\operatorname{Re}(s) > 1/2$. The collision
fluctuation sum converges for $\operatorname{Re}(s) > 2$ and
appears to diverge logarithmically at $s = 1$.
Computation suggests that $s = 1$ is the boundary between
convergence and divergence for $\Phi(s, \ell)$, with
Mertens-type logarithmic growth rather than PNT-type
convergence.

\subsection*{Forward direction}

The divergence at $s = 1$ arises from a persistent per-class
bias in $\Delta_p$. The next paper investigates whether
subtracting this bias (centering) produces a sum that converges
at $s = 1$.

% ============================================================
\begin{thebibliography}{9}

\bibitem{paper1}
A.~S.~Petty,
Why the golden ratio selects the prime three,
research note, 2020.

\bibitem{paper11}
A.~S.~Petty,
Bin derangements and the gate width theorem,
research note, 2024.

\bibitem{paper12}
A.~S.~Petty,
Silent primes and the variance of the collision count,
research note, 2024.

\bibitem{mertens}
F.~Mertens,
Ein Beitrag zur analytischen Zahlentheorie,
\emph{J.\ Reine Angew.\ Math.}\ \textbf{78} (1874), 46--62.

\bibitem{davenport}
H.~Davenport,
\emph{Multiplicative Number Theory},
3rd~ed., Springer, 2000.

\bibitem{hardy}
G.~H.~Hardy and E.~M.~Wright,
\emph{An Introduction to the Theory of Numbers},
6th~ed., Oxford University Press, 2008.

\end{thebibliography}

\end{document}
